\section{مدل}

\begin{frame}{چه ساخته شد و مدل به چه چیزی نگاه می‌کند}
  \begin{itemize}
    \item ۱۰۰ ویژگی، ۱۳ خانواده‌ی مدل، غربال چهارمرحله‌ای، انتخاب نهایی یک مدل درختی
      کوانتایلی
    \item پروتکل: تقسیم زمانی با مبدأ غلتان؛ پنجره‌ی آزمون فقط یک‌بار باز شد
    \item مهم‌ترین ورودی‌ها: هویت سلف، برهم‌کنش روزهفته با شهر، حجم رزرو، نرخ دیروزِ همان
      سلول، شوک دیروز
  \end{itemize}
  \begin{figure}
    \centering
    \includegraphics[width=0.62\textwidth]{9.1_perm_importance_fa}
  \end{figure}
  \note{کلمه‌ی SHAP را نگو، بگو «سهم هر ویژگی در تصمیم مدل». نکته‌ی مهم اینکه این فهرست
  دقیقاً با یافته‌های بخش ۳ می‌خواند، یعنی مدل چیز عجیبی یاد نگرفته. شاهد:
  \lr{reports/phase9/9.2\_shap\_mean\_abs.csv}، \lr{reports/phase7/model\_comparison.md}.}
\end{frame}

\begin{frame}{نتایج منفی}
  \begin{itemize}
    \item از ۲۵ مدل فقط ۵ تا برد آماری معنادار داشتند
    \item شبکه‌های عصبی و مدل سطح فرد کمکی نکردند
    \item قهرمان اعتبارسنجی و قهرمان پنجره‌ی آزمون یکی نبودند، و طبق پروتکل قهرمان عوض نشد
  \end{itemize}
  \begin{figure}
    \centering
    \includegraphics[width=0.62\textwidth]{slides_15_funnel}
  \end{figure}
  \note{این اسلاید را حذف نکن. استادها روی همین قضاوت می‌کنند. شاهد: یافته‌های ۲۸، ۳۱، ۳۶؛
  ردیف ۴۹ decision\_log.}
\end{frame}
